\section*{Capítulo \ref{ch:g6:rocks-and-minerals} --- Rocas y minerales}

\begin{solution}{exo:g6:rocks-and-minerals:1}
Granito: ígnea. Arenisca: sedimentaria. Mármol: metamórfica. Basalto:
ígnea. Caliza: sedimentaria. Pizarra: metamórfica.
\end{solution}

\begin{solution}{exo:g6:rocks-and-minerals:2}
El cuarzo es un \omterm{def:g6:rocks-and-minerals:mineral}{mineral}; el granito es una \omterm{def:g3:stones-pebbles-sand:rock}{roca}. Un \omterm{def:g6:rocks-and-minerals:mineral}{mineral} es una sola
sustancia natural con sus átomos ordenados de forma regular; una \omterm{def:g3:stones-pebbles-sand:rock}{roca} está
hecha de uno o de varios \omterm{def:g6:rocks-and-minerals:mineral}{minerales} unidos entre sí (el granito: feldespato,
cuarzo y mica).
\end{solution}

\begin{solution}{exo:g6:rocks-and-minerals:3}
Talco (1), calcita (3), cuarzo (7), diamante (10). El diamante, el más
duro, raya a todos los demás.
\end{solution}

\begin{solution}{exo:g6:rocks-and-minerals:4}
``La \omterm{def:g3:stones-pebbles-sand:rock}{roca} fundida se convierte en una \omterm{def:g6:rocks-and-minerals:rock-families}{roca ígnea} por enfriamiento. Una
caliza se convierte en un mármol por el calor y la presión. Un granito se
convierte en granos de \omterm{def:g3:stones-pebbles-sand:sand}{arena} por \omterm{def:g5:water-shapes-the-land:weathering}{meteorización}''.
\end{solution}

\begin{solution}{exo:g6:rocks-and-minerals:5}
(a) Basalto. (b) Mármol. (c) Arenisca.
\end{solution}

\begin{solution}{exo:g6:rocks-and-minerals:6}
Su \omterm{def:g6:rocks-and-minerals:hardness}{dureza} es mayor que la de una uña (unos 2.5) y menor que la de una
moneda de cobre (unos 3.5). Podría ser la calcita (3).
\end{solution}

\begin{solution}{exo:g6:rocks-and-minerals:7}
Cuarzo: $21 \div 60 = 0.35$, es decir, el 35~\%. Feldespato:
$33 \div 60 = 0.55$, es decir, el 55~\%. Mica: $6 \div 60 = 0.10$, es
decir, el 10~\%. Comprobación: $35 + 55 + 10 = 100$.
\end{solution}

\begin{solution}{exo:g6:rocks-and-minerals:8}
El granito se enfría muy despacio, en las profundidades: sus granos tienen
mucho tiempo para crecer y se hacen grandes. El basalto viene de la \omterm{def:g4:volcanoes-earthquakes-first-look:eruption}{lava},
que se enfría deprisa en la superficie: sus granos no tienen tiempo de
crecer y se quedan demasiado pequeños para verlos.
\end{solution}

\begin{solution}{exo:g6:rocks-and-minerals:9}
(a) La caliza es sedimentaria; el mármol es metamórfico. (b) En las
profundidades, el calor y la presión la transformaron sin fundirla: su
calcita creció en forma de cristales brillantes que encajan entre sí.
(c) Las dos están hechas del mismo \omterm{def:g6:rocks-and-minerals:mineral}{mineral}, la calcita, y es la calcita la
que burbujea con el ácido.
\end{solution}

\begin{solution}{exo:g6:rocks-and-minerals:10}
(a) Sí: el cuarzo (7) es más duro que el vidrio (unos 5.5). (b) No: la
fluorita (4) es más blanda que el vidrio. (c) Sí: un clavo de acero (unos
6.5) es más duro que el feldespato (6).
\end{solution}

\begin{solution}{exo:g6:rocks-and-minerals:11}
Por ejemplo: (1) el granito de la \omterm{def:g1:landscapes-around-us:landform}{montaña} (ígnea) se rompe por
\omterm{def:g5:water-shapes-the-land:weathering}{meteorización}, y sus granos de cuarzo quedan sueltos; (2) un \omterm{def:g1:landscapes-around-us:water}{río} se los
lleva: \omterm{def:g5:water-shapes-the-land:erosion}{erosión}; (3) los granos se depositan en capas en el fondo del \omterm{def:g2:water-on-earth:ocean}{mar}:
\omterm{def:g5:water-shapes-the-land:deposition}{sedimentación}; (4) enterrados bajo otras capas, se comprimen y se cementan
formando una arenisca: endurecimiento (la \omterm{def:g3:stones-pebbles-sand:rock}{roca} es ahora sedimentaria);
(5) llevada a las profundidades, la arenisca se transforma por el calor y
la presión en una \omterm{def:g6:rocks-and-minerals:rock-families}{roca metamórfica}.
\end{solution}

\begin{solution}{exo:g6:rocks-and-minerals:12}
(a) No se raya con el acero (unos 6.5) y raya el vidrio: es más duro que
6.5. Lo raya el topacio (8): es más blando que 8. Su \omterm{def:g6:rocks-and-minerals:hardness}{dureza} es 7: es el
cuarzo. (b) El cuarzo es más duro que casi todo lo que roza con él (otros
granos, \omterm{def:g3:stones-pebbles-sand:pebble}{cantos}): apenas se desgasta. La calcita (3) y la mica (de 2 a 3) son
mucho más blandas: el roce las muele en trocitos diminutos que el agua se
lleva.
\end{solution}

\begin{solution}{pb:g6:rocks-and-minerals:1}
\textbf{1.} Granito, una \omterm{def:g6:rocks-and-minerals:rock-families}{roca ígnea}; sus \omterm{def:g6:rocks-and-minerals:mineral}{minerales} son el feldespato (rosa
o blanco), el cuarzo (gris, vidrioso) y la mica (negra, brillante).

\textbf{2.} Caliza, una \omterm{def:g6:rocks-and-minerals:rock-families}{roca sedimentaria}. Las conchas eran de animales que
vivían en el \omterm{def:g2:water-on-earth:ocean}{mar}; sus restos se depositaron y se endurecieron hasta formar
la \omterm{def:g3:stones-pebbles-sand:rock}{roca}.

\textbf{3.} Pizarra, una \omterm{def:g6:rocks-and-minerals:rock-families}{roca metamórfica}, formada a partir de la lutita.

\textbf{4.} Mármol, una \omterm{def:g6:rocks-and-minerals:rock-families}{roca metamórfica}, formado a partir de la caliza.

\textbf{5.} Cuarzo: no lo raya el acero (unos 6.5) y raya el vidrio (unos
5.5); su \omterm{def:g6:rocks-and-minerals:hardness}{dureza} es 7.

\textbf{6.} Los raya el acero (unos 6.5) pero no el vidrio (unos 5.5):
entre 5.5 y 6.5. Es el feldespato (ortoclasa, 6).

\textbf{7.} La mica (\omterm{def:g6:rocks-and-minerals:hardness}{dureza} de 2 a 3).

\textbf{8.} Sí: el pavimento es de caliza, hecha de calcita (3), más blanda
que una moneda de cobre (unos 3.5).

\textbf{9.} La \omterm{def:g5:water-shapes-the-land:weathering}{meteorización}: la \omterm{def:g4:clouds-rain-snow:precipitation}{lluvia} desgasta lentamente la superficie
de la piedra.

\textbf{10.} La \omterm{def:g1:landscapes-around-us:landform}{montaña} se meteoriza; sus granos de cuarzo son llevados por
los \omterm{def:g1:landscapes-around-us:water}{ríos} y se depositan en capas, en el fondo de un \omterm{def:g1:landscapes-around-us:water}{lago} o del \omterm{def:g2:water-on-earth:ocean}{mar};
enterrados y endurecidos, pueden convertirse en una arenisca, una \omterm{def:g6:rocks-and-minerals:rock-families}{roca
sedimentaria}.

\textbf{11.} Las tres familias están presentes: ígnea (el granito de la
fuente), sedimentaria (la caliza del pavimento) y metamórfica (la pizarra
del tejado y el mármol de la estatua).

\textbf{12.} El \omterm{def:g6:rocks-and-minerals:mineral}{mineral} más duro es el cuarzo del granito: \omterm{def:g6:rocks-and-minerals:hardness}{dureza} de Mohs
\textbf{7}.
\end{solution}
